\documentclass[11pt,letterpaper]{article}
\usepackage[margin=1in]{geometry}
\usepackage{amsmath,amssymb,amsthm}
\usepackage{physics}
\usepackage{hyperref}
\usepackage{cleveref}
\usepackage{booktabs}

\numberwithin{equation}{section}
\renewcommand{\theequation}{S16.\arabic{equation}}

\title{Supplement S16: QCD as Angular Geometry and the \\
Derived Nuclear Force \\
\large Supporting Manuscript \S19}
\author{UDT Collaboration}
\date{}

\begin{document}
\maketitle

\begin{abstract}
The 8~tensor operators on the $\ell = 1$ subspace of $S^2$ close as
the Lie algebra $\mathfrak{su}(3)$ with Casimir $C_2 = (4/3)I_3$,
verified to $10^{-16}$.  This renders the gauge principle scaffolding.
Four nuclear coupling constants --- $g_A = 4/\pi$ ($-0.23\%$),
$f_\pi = 180\,m_e$ ($-0.25\%$), $g_{\pi NN} = 2\pi^4/15$ ($+0.06\%$),
$g^2/(4\pi) = \pi^7/225$ ($-0.6\%$) --- follow from
$(j, \ell, |\kappa_\mathrm{max}|) = (1/2, 1, 3)$ and $m_e$.  The
one-pion exchange potential gives $V(1.5~\text{fm}) = -41$~MeV; the
scalar $\phi$ sector contributes zero ($\alpha_\mathrm{phys} \sim
10^{-40}$).  The proton charge radius $R_p = \hbar c/(C\varphi_\mathrm{gold})
= 0.865$~fm ($-1.1\%$) sits between the muonic and electronic puzzle
values.  Verified by \texttt{analysis/22\_su3\_closure.py} and
\texttt{analysis/23\_nuclear\_couplings.py}.
\end{abstract}

\tableofcontents

%----------------------------------------------------------------------
\section{The $\mathfrak{su}(3)$ algebra from $S^2$}
\label{sec:su3}
%----------------------------------------------------------------------

The Dirac equation on the metric decomposes angularly into spinor
harmonics $\Omega_\kappa^{m_j}(\theta, \varphi)$.  At $\ell = 1$,
the angular space is 3-dimensional, supporting:
\begin{itemize}
  \item 3~rank-1 operators (angular momentum components $L_i$),
  \item 5~rank-2 operators (quadrupole-type couplings),
\end{itemize}
for a total of $8 = (2\ell+1)^2 - 1$ traceless Hermitian operators.

\begin{theorem}[su(3) closure]
The 8~traceless Hermitian operators on the $\ell = 1$ subspace satisfy
\begin{equation}
  [T_a, T_b] = i\,f_{abc}\,T_c\,,
  \qquad C_2 = \sum_{a=1}^8 T_a^2 = \frac{4}{3}\,I_3\,,
\end{equation}
with the standard $\mathfrak{su}(3)$ structure constants $f_{abc}$.
\end{theorem}

This is verified numerically to residual $< 1.1 \times 10^{-16}$
(machine precision).

\subsection{Geometric protection}

The angular metric $g_{\theta\theta} = r^2$ carries no $\phi$
dependence.  The $\mathfrak{su}(3)$ algebra is therefore
\emph{geometrically protected} from the dilation field: it closes
identically regardless of $\phi(r)$.  This is why QCD is
asymptotically free --- at short distances, the angular algebra
dominates; the radial $\phi$ dynamics cannot modify it.

\subsection{The gauge principle as scaffolding}

The Standard Model \emph{postulates} $\mathrm{SU}(3) \times
\mathrm{SU}(2) \times \mathrm{U}(1)$ as the gauge group and constructs
the theory from gauge invariance.  In UDT, the same group structure
arises from the metric's angular decomposition:
\begin{align}
  \mathrm{SU}(3) &: \text{rank-1 + rank-2 on } \ell = 1\,,\\
  \mathrm{SU}(2) &: \text{rank-1 on } j = 1/2\,,\\
  \mathrm{U}(1) &: \text{Coulomb from } g_{tt}g_{rr} = -c^2\,.
\end{align}
The gauge principle is not wrong --- it is redundant.  The metric
provides everything the gauge postulate was designed to ensure.

%----------------------------------------------------------------------
\section{Nuclear coupling constants}
\label{sec:nuclear}
%----------------------------------------------------------------------

Four nuclear quantities form a closed algebraic chain:

\subsection{Axial coupling $g_A$}

\begin{equation}\label{eq:gA}
  g_A = \frac{(2j+1)^2}{\pi} = \frac{4}{\pi} = 1.2732
  \qquad(\text{exp: }1.2762,\ {-0.23\%})\,.
\end{equation}
The numerator $(2j+1)^2 = 4$ is the spin bilinear dimension.  The
denominator $\pi$ is the angular circumference.

\textbf{Deep derivation:}  Including the SU(6) quark model factor
$(N_c + 2)/3 = 5/3$:
\begin{equation}
  g_A = \frac{(N_c+2)(2j+1)^2}{(2|\kappa_\mathrm{max}|-1)\pi}
  = \frac{5 \times 4}{5\pi} = \frac{4}{\pi}\,.
\end{equation}
The cancellation $(N_c + 2) = (2\ell + 3) = (2|\kappa_\mathrm{max}|-1) = 5$
is specific to the Diophantine triple $(1/2, 1, 3)$.

\subsection{Pion decay constant $f_\pi$}

\begin{equation}\label{eq:fpi}
  f_\pi = \binom{9}{2}(2|\kappa_\mathrm{max}|-1)\,m_e
  = 36 \times 5 \times m_e = 180\,m_e = 91.98~\text{MeV}\,,
\end{equation}
compared to $92.21$~MeV ($-0.25\%$).  The factor $36 = \binom{9}{2}$
is the same antisymmetric coupling count as in $1/\alpha_\mathrm{EM}$;
$5 = 2|\kappa_\mathrm{max}|-1$ is the angular content.

The ratio $f_\pi/m_\pi$ is algebraic:
\begin{equation}
  \frac{f_\pi}{m_\pi}
  = \frac{(2\ell+1)(2|\kappa_\mathrm{max}|-1)}{(2|\kappa_\mathrm{max}|+1)\pi}
  = \frac{15}{7\pi}\,.
\end{equation}

\subsection{Pion--nucleon coupling (Goldberger--Treiman)}

The Goldberger--Treiman relation with all derived inputs:
\begin{equation}\label{eq:gpiNN}
  g_{\pi NN} = \frac{g_A\,m_N}{f_\pi}
  = \frac{(4/\pi) \times 6\pi^5\,m_e}{180\,m_e}
  = \frac{24\pi^4}{180} = \frac{2\pi^4}{15} = 12.99\,,
\end{equation}
compared to $\sim\!12.98$ ($+0.06\%$).  The relation is used with
\emph{physical} (not bare) quantities --- no renormalization bridge
is needed.

\subsection{Nuclear coupling constant}

\begin{equation}\label{eq:g2}
  \frac{g^2_{\pi NN}}{4\pi}
  = \frac{(2\pi^4/15)^2}{4\pi}
  = \frac{\pi^7}{225} = 13.42\,,
\end{equation}
compared to $\sim\!13.5$ ($-0.6\%$).  The power $\pi^7$ has exponent
$7 = 2|\kappa_\mathrm{max}| + 1$ (the closure number, the same 7 as
CMB peaks).  The denominator $225 = 15^2 = [(2\ell+1)(2|\kappa_\mathrm{max}|-1)]^2$.

%----------------------------------------------------------------------
\section{The one-pion exchange potential}
\label{sec:opep}
%----------------------------------------------------------------------

The OPEP with derived coupling gives the deuteron-channel (${}^3S_1$)
central potential:
\begin{equation}
  V_\mathrm{OPEP}(r) = -\frac{g^2_{\pi NN}}{4\pi}
  \frac{m_\pi^2}{12\,m_N^2}\,\frac{e^{-m_\pi r/(\hbar c)}}{r/(\hbar c)}\,.
\end{equation}

\begin{center}
\begin{tabular}{rc}
\toprule
$r$ [fm] & $V$ [MeV] \\
\midrule
1.0 & $-87$ \\
1.5 & $-41$ \\
2.0 & $-22$ \\
3.0 & $-7.4$ \\
5.0 & $-1.1$ \\
\bottomrule
\end{tabular}
\end{center}

Adding $\omega$-meson exchange ($m_\omega = 784.4$~MeV, the
$\kappa = -1$, $n = 2$ eigenvalue at $+0.22\%$) produces the full
shape: repulsive core ($r < 0.7$~fm), attractive well
($0.7 < r < 3$~fm), and Yukawa tail ($r > 3$~fm).

\subsection{Scalar sector: zero contribution}

The scalar $\phi$ backreaction through $G_N$ gives
$\alpha_\mathrm{phys} \sim 10^{-40}$.  The nuclear force comes
\emph{entirely} from the angular sector (pion exchange), consistent
with the identification of the nuclear force as a QCD residual.

%----------------------------------------------------------------------
\section{Proton charge radius}
\label{sec:rp}
%----------------------------------------------------------------------

The charge radius is the geometric ratio of baryon to meson cavity:
\begin{equation}\label{eq:rp}
  R_p = \frac{\hbar c}{C\,\varphi_\mathrm{gold}}
  = \frac{r_{*b}}{r_*} \times \frac{\hbar c}{C}
  = 0.865~\text{fm}\,,
\end{equation}
compared to CODATA $0.8751$~fm ($-1.1\%$).  This sits between the
muonic hydrogen value ($0.841$~fm) and the electronic value
($0.877$~fm).

%----------------------------------------------------------------------
\section{Verification}
%----------------------------------------------------------------------

The $\mathfrak{su}(3)$ closure is verified by
\texttt{analysis/22\_su3\_closure.py}; all nuclear couplings by
\texttt{analysis/23\_nuclear\_couplings.py}.

Outputs: \texttt{data/generated/22\_su3\_closure.json},
\texttt{data/generated/23\_nuclear\_couplings.json}.

\end{document}
